\documentclass[12pt]{article}
\usepackage[a4paper,margin=1.25cm]{geometry}
\usepackage{amsmath,amssymb,mathtools}
\usepackage{graphicx}
\usepackage{booktabs}
\usepackage{tabularx}
\usepackage{tikz}
\usepackage[T1]{fontenc}
\usepackage[utf8]{inputenc}
\usepackage[english]{babel}

\setlength{\parindent}{0pt}
\setlength{\parskip}{0.45em}
\setlength{\abovedisplayskip}{5pt plus 1pt minus 1pt}
\setlength{\belowdisplayskip}{5pt plus 1pt minus 1pt}
\setlength{\abovedisplayshortskip}{3pt plus 1pt minus 1pt}
\setlength{\belowdisplayshortskip}{3pt plus 1pt minus 1pt}

\newcommand{\dd}{\,\mathrm{d}}
\newcommand{\e}{\mathrm{e}}
\newcommand{\An}{A_n}
\newcommand{\Bn}{B_n}

\begin{document}

\begin{center}
{\Large\bfseries A logarithmic convolution}\\[0.35em]
{\large almost every index has \(\log k\approx\log n\)}
\end{center}

\[
\boxed{\displaystyle
\lim_{n\to\infty}
\frac{\log^2 n}{n}\sum_{k=2}^{n-2}\frac1{\log k\,\log(n-k)}
=1 .}
\]

\textbf{Overview.}
Put
\[
x_k=\frac{\log n}{\log k},\qquad 2\le k\le n-2,
\qquad\text{so that}\qquad
\An:=\frac{\log^2n}{n}\sum_{k=2}^{n-2}\frac1{\log k\,\log(n-k)}
=\frac1n\sum_{k=2}^{n-2}x_kx_{n-k}.
\]
Every \(x_k\ge1\), which gives the lower bound at once. For the upper
bound, Cauchy--Schwarz plus the reflection \(k\mapsto n-k\) replaces the
product \(x_kx_{n-k}\) by \(x_k^2\), and then the key structural fact
takes over: for any \(\alpha<1\), all but \(n^{\alpha}\) of the indices
satisfy \(k\ge n^{\alpha}\), hence \(\log k\ge\alpha\log n\) and
\(x_k\le\frac1\alpha\). The exceptional indices are so few, a power of
\(n\) below \(n\), that even the worst possible value
\(x_2=\frac{\log n}{\log2}\) cannot rescue them.

In words: the logarithm is so slowly varying that \emph{almost every}
\(k\le n\) has \(\log k\) essentially equal to \(\log n\), so the sum has
\(n\) terms each essentially equal to \(\frac1{\log^2n}\).

\section*{1. Notation and the lower bound}

For \(2\le k\le n-2\) both \(k\le n\) and \(n-k\le n\), so
\(\log k\le\log n\) and \(\log(n-k)\le\log n\), whence
\[
x_k=\frac{\log n}{\log k}\ge1,
\qquad
x_{n-k}=\frac{\log n}{\log(n-k)}\ge1 .
\tag{1}
\]
Therefore each of the \(n-3\) summands of
\[
\An=\frac1n\sum_{k=2}^{n-2}x_kx_{n-k}
\tag{2}
\]
is at least \(1\), and
\[
\boxed{\displaystyle \An\ge\frac{n-3}{n}.}
\tag{3}
\]
In particular \(\liminf_n\An\ge1\). The whole problem is the matching
upper bound.

\section*{2. Removing the convolution by Cauchy--Schwarz}

Apply the Cauchy--Schwarz inequality to (2):
\[
\sum_{k=2}^{n-2}x_kx_{n-k}
\le\sqrt{\left(\sum_{k=2}^{n-2}x_k^2\right)
\left(\sum_{k=2}^{n-2}x_{n-k}^2\right)} .
\]
The substitution \(k\mapsto n-k\) is a bijection of the index range
\(\{2,\dots,n-2\}\) onto itself, so the two factors are equal and
\[
\An\le\Bn:=\frac1n\sum_{k=2}^{n-2}x_k^2
=\frac{\log^2n}{n}\sum_{k=2}^{n-2}\frac1{\log^2k}.
\tag{4}
\]
By the same reasoning as (1), also \(\Bn\ge\frac{n-3}{n}\). The
convolution structure has been removed at no cost: it remains to show
\(\Bn\to1\).

\section*{3. The upper bound: splitting at \(n^{\alpha}\)}

Fix \(\alpha\in(0,1)\) and split the index range at \(n^{\alpha}\):
\[
\Bn=\frac1n\sum_{2\le k<n^{\alpha}}x_k^2+\frac1n\sum_{n^{\alpha}\le k\le n-2}x_k^2 .
\tag{5}
\]

\textbf{Small indices.} There are fewer than \(n^{\alpha}\) of them, and
the largest possible value of \(x_k\) on the whole range is
\(x_2=\frac{\log n}{\log2}\). Hence
\[
\frac1n\sum_{2\le k<n^{\alpha}}x_k^2
\le\frac{n^{\alpha}}{n}\left(\frac{\log n}{\log2}\right)^2
=n^{\alpha-1}\left(\frac{\log n}{\log2}\right)^2
\xrightarrow[n\to\infty]{}0,
\tag{6}
\]
since \(\alpha-1<0\) and a power beats a square of a logarithm.

\textbf{Large indices.} If \(k\ge n^{\alpha}\) then
\(\log k\ge\alpha\log n\), so \(x_k\le\frac1\alpha\) and
\[
\frac1n\sum_{n^{\alpha}\le k\le n-2}x_k^2
\le\frac1n\cdot n\cdot\frac1{\alpha^2}
=\frac1{\alpha^2}.
\tag{7}
\]

Now let \(\varepsilon>0\). Choose \(\alpha<1\) close enough to \(1\) that
\[
\frac1{\alpha^2}<1+\frac\varepsilon2 ,
\]
and then \(N\) so large that for \(n\ge N\) both
\[
n^{\alpha-1}\left(\frac{\log n}{\log2}\right)^2<\frac\varepsilon2
\qquad\text{and}\qquad
\frac{n-3}{n}>1-\varepsilon .
\]
Combining (3), (4), (5), (6), (7), for every \(n\ge N\)
\[
1-\varepsilon<\frac{n-3}{n}\le\An\le\Bn
\le\frac\varepsilon2+\frac1{\alpha^2}<1+\varepsilon .
\]
Since \(\varepsilon\) was arbitrary,
\[
\boxed{\displaystyle
\lim_{n\to\infty}\Bn=1
\qquad\text{and}\qquad
\lim_{n\to\infty}\An=1 .}
\]

Note the order of the choices: \(\alpha\) is selected from
\(\varepsilon\) first, and only then is \(N\) selected from both. The
exponent \(\alpha\) is a free parameter that trades the two error terms
against each other: pushing \(\alpha\) towards \(1\) improves (7) and
worsens (6), but (6) has a power of \(n\) to spend and always wins.

\section*{4. How slow the convergence is}

The proof gives no useful rate, and for good reason: the true rate is
logarithmic.

\begin{center}
\small
\renewcommand{\arraystretch}{1.35}
\begin{tabularx}{0.9\textwidth}{@{}lXXX@{}}
\toprule
\(n\) & \(\An\) & \(\An-1\) & \(2/\log n\) \\
\midrule
\(10^2\) & \(1.7792\) & \(0.7792\) & \(0.4343\) \\
\(10^3\) & \(1.4603\) & \(0.4603\) & \(0.2895\) \\
\(10^4\) & \(1.3002\) & \(0.3002\) & \(0.2171\) \\
\(10^5\) & \(1.2204\) & \(0.2204\) & \(0.1737\) \\
\bottomrule
\end{tabularx}
\end{center}

At \(n=10^5\) the sequence is still \(22\%\) above its limit. The reason
is visible from a continuous model. Write \(k=tn\) with \(t\in(0,1)\) and
\(L=\log n\); then
\[
\log k=L+\log t,
\qquad
x_k=\frac{L}{L+\log t}=\frac1{1+\frac{\log t}{L}}
\approx1-\frac{\log t}{L},
\]
and similarly for \(x_{n-k}\) with \(t\) replaced by \(1-t\). Hence
\[
x_kx_{n-k}\approx1-\frac{\log t+\log(1-t)}{L},
\]
and averaging over \(t\in(0,1)\), using
\(\int_0^1(-\log t)\dd t=\int_0^1(-\log(1-t))\dd t=1\),
\[
\boxed{\displaystyle \An\approx1+\frac2{\log n}.}
\tag{8}
\]
The last column of the table is this prediction. It captures the right
order of magnitude and the right constant asymptotically, though at these
values of \(n\) the neglected \(O(\log^{-2}n)\) term is still comparable
to the main correction, which is what one expects when the
expansion parameter is \(1/\log n\approx0.09\) at \(n=10^5\).

\section*{5. What made it work}

The only property of \(\log\) used in Section 3 was
\[
\log(n^{\alpha})=\alpha\log n ,
\tag{9}
\]
i.e. that the logarithm of a fixed power of \(n\) is a fixed fraction of
\(\log n\). This is what forces \(x_k\) to be uniformly bounded on all
but a negligible set of indices, and it is a statement about slow
variation, not about the logarithm specifically.

\begin{center}
\small
\renewcommand{\arraystretch}{1.45}
\begin{tabularx}{\textwidth}{@{}>{\raggedright\arraybackslash}p{0.28\textwidth}X@{}}
\toprule
Weight \(w(k)\) & Behaviour of \(\dfrac{w(n)^2}{n}\sum_{k}\dfrac1{w(k)w(n-k)}\) \\
\midrule
\(w(k)=\log k\)
& Limit \(1\), by the argument above. \\
\addlinespace
\(w(k)=\log^{p}k\), \(p>0\)
& Limit \(1\): \(\bigl(\frac{\log n}{\log k}\bigr)^{p}\le\alpha^{-p}\) on
\(k\ge n^{\alpha}\), and the small-index block is still killed by
\(n^{\alpha-1}\log^{2p}n\to0\). \\
\addlinespace
\(w(k)=k^{\beta}\), \(\beta>0\)
& \emph{Not} \(1\): here \(\frac{w(n)^2}{w(k)w(n-k)}=\bigl(t(1-t)\bigr)^{-\beta}\)
does not concentrate, and the limit is the Beta integral
\(\int_0^1 t^{-\beta}(1-t)^{-\beta}\dd t=B(1-\beta,1-\beta)\)
for \(0<\beta<1\), diverging for \(\beta\ge1\). \\
\bottomrule
\end{tabularx}
\end{center}

The contrast in the last row is the point of the problem. For a power
weight the rescaled summand is a genuine function of \(t=k/n\) and the
limit is a Riemann integral, generally not \(1\). For a logarithmic
weight the rescaled summand converges to the constant \(1\) for
\emph{every} fixed \(t\in(0,1)\), and the limit is \(\int_0^1 1\dd t=1\).
Formula (8) is the first correction to that degeneracy, and it is the
Riemann integral of \(-\log t-\log(1-t)\), divided by \(\log n\).

\end{document}
